% !TEX root = ../main.tex
% =====================================================================
%  CHƯƠNG 3 — THIẾT KẾ MẠCH ĐỌC LOAD CELL, CHUYỂN ĐỔI ÁP THÀNH DÒNG
%             VÀ MẠCH ĐỌC DÒNG
% ---------------------------------------------------------------------
%  Người phụ trách : [TODO: Họ tên - MSSV]
%  Tiền tố label   : c3      (fig:c3:..., tab:c3:..., eq:c3:..., sec:c3:...)
%  Thư mục ảnh     : image/chuong3/
%  File TLTK       : references/chuong3.bib   (key: c3_...)
%  Độ dài gợi ý    : 10–15 trang
%  Đáp ứng yêu cầu : (a), (b) và phần mạch đọc dòng của (e).
%  Quy trình MỖI KHỐI: (1) sơ đồ mạch -> (2) tính toán, phương trình đánh số
%                    -> (3) chọn giá trị linh kiện -> (4) MÔ PHỎNG RIÊNG LẺ
%                    -> (5) \nhanxet{} so sánh với lý thuyết.
%  Giao tiếp       : Dải tín hiệu vào/ra phải khớp Bảng tab:c2:giao-tiep-khoi.
% =====================================================================
\chapter{Thiết kế mạch đọc Load cell, chuyển đổi áp thành dòng và mạch đọc dòng}
\label{chap:c3}

% ---------------------------------------------------------------------
\section{Yêu cầu thiết kế}
\label{sec:c3:yeu-cau}

\TODO{Tóm tắt yêu cầu của các khối trong chương: dải vào/ra theo Bảng~\ref{tab:c2:giao-tiep-khoi}, nguồn cấp, sai số mục tiêu.}

% ---------------------------------------------------------------------
\section{Lựa chọn và mô phỏng Load cell}
\label{sec:c3:loadcell}
% Yêu cầu (a).

\subsection{Lựa chọn Load cell}
\label{sec:c3:chon-loadcell}

\TODO{Lý do lựa chọn Load cell, bảng thông số trích từ datasheet (có \textbackslash cite).}

\subsection{Mô hình hóa Load cell}
\label{sec:c3:mo-hinh-loadcell}

\TODO{Xây dựng mô hình điện trở tương đương, tính $R_f$ theo khối lượng $m$.}

\begin{equation}
    \left\{
    \begin{aligned}
        R_{\mathrm{in}}  &= \TODO{} \\
        R_{\mathrm{out}} &= \TODO{}
    \end{aligned}
    \right.
    \label{eq:c3:rin-rout}
\end{equation}

\subsection{Mô phỏng Load cell}
\label{sec:c3:mo-phong-loadcell}

\haihinh{chuong3/loadcell_m0.png}{Mô phỏng Load cell khi $m = \qty{0}{\kilo\gram}$}{fig:c3:loadcell-m0}%
        {chuong3/loadcell_m1.png}{Mô phỏng Load cell khi $m = \qty{1}{\kilo\gram}$}{fig:c3:loadcell-m1}

\nhanxet{\TODO{So sánh kết quả ở Hình~\ref{fig:c3:loadcell-m0} và
Hình~\ref{fig:c3:loadcell-m1} với giá trị tính theo \eqref{eq:c3:rin-rout}.}}

% ---------------------------------------------------------------------
\section{Lựa chọn và mô phỏng OPAMP theo datasheet}
\label{sec:c3:opamp}
% Yêu cầu (b). Chọn INA cho khuếch đại vi sai và OPAMP cho V-I, đọc dòng, đệm.
% Mô phỏng các thông số (Vos, Ib, GBW, CMRR...) và so sánh với datasheet.

Bộ khuếch đại dụng cụ được lựa chọn dựa trên các thông số trong datasheet~\cite{c3_ina333}.

\TODO{Lý do lựa chọn INA / OPAMP.}

\begin{table}[H]
    \centering
    \caption{So sánh thông số OPAMP giữa datasheet và mô phỏng}
    \label{tab:c3:opamp-datasheet}
    \renewcommand{\arraystretch}{1.3}
    \begin{tabular}{|l|c|c|}
        \hline
        \textbf{Thông số} & \textbf{Datasheet} & \textbf{Mô phỏng} \\ \hline
        Điện áp offset ngõ vào $V_{\mathrm{OS}}$ & \TODO{} & \TODO{} \\ \hline
        Dòng phân cực $I_{\mathrm{B}}$           & \TODO{} & \TODO{} \\ \hline
        Tích độ lợi -- băng thông (GBW)          & \TODO{} & \TODO{} \\ \hline
        CMRR                                     & \TODO{} & \TODO{} \\ \hline
    \end{tabular}
\end{table}

% ---------------------------------------------------------------------
\section{Mạch khuếch đại vi sai}
\label{sec:c3:khuech-dai}

\TODO{Tính độ lợi cần thiết theo Bảng~\ref{tab:c2:giao-tiep-khoi}, chọn $R_G$, mạch hiệu chỉnh điểm không.}

\subsection{Mô phỏng mạch khuếch đại vi sai}
\label{sec:c3:mo-phong-khuech-dai}

\TODO{Hình mô phỏng tại $m = \qty{0}{\kilo\gram}$ và $m = \qty{1}{\kilo\gram}$, kèm \textbackslash nhanxet.}

% ---------------------------------------------------------------------
\section{Mạch chuyển đổi áp thành dòng 4--20mA}
\label{sec:c3:vi-converter}

\TODO{Thiết lập quan hệ tuyến tính $I = aV + b$, thiết kế mạch, chọn linh kiện.}

\subsection{Mô phỏng mạch chuyển đổi áp thành dòng}
\label{sec:c3:mo-phong-vi}

\TODO{Hình mô phỏng tại các mức điện áp vào tương ứng \qty{0}{\kilo\gram} và \qty{1}{\kilo\gram}, kèm \textbackslash nhanxet.}

% ---------------------------------------------------------------------
\section{Mạch đọc dòng}
\label{sec:c3:doc-dong}
% Phần mạch đọc dòng của yêu cầu (e). Điện trở shunt + khuếch đại vi sai + đệm.
% Ngõ ra phải nằm trong dải ngõ vào của ADC (Chương 5).

\TODO{Thiết kế mạch đọc dòng, tính $\Rshunt$ và độ lợi.}

\subsection{Mô phỏng mạch đọc dòng}
\label{sec:c3:mo-phong-doc-dong}

\TODO{Hình mô phỏng tại \qty{4}{\milli\ampere} và \qty{20}{\milli\ampere}, kèm \textbackslash nhanxet.}

% ---------------------------------------------------------------------
\section{Tổng kết kết quả mô phỏng các khối}
\label{sec:c3:tong-ket}

\begin{table}[H]
    \centering
    \caption{Tổng hợp kết quả mô phỏng riêng lẻ các khối ở Chương~\ref{chap:c3}}
    \label{tab:c3:tong-ket}
    \renewcommand{\arraystretch}{1.3}
    \begin{tabularx}{\textwidth}{|X|c|c|c|c|}
        \hline
        \textbf{Khối / Đại lượng} & \textbf{Điều kiện} & \textbf{Lý thuyết} & \textbf{Mô phỏng} & \textbf{Sai số (\unit{\percent})} \\ \hline
        Load cell, $\Vout$          & $m = \qty{1}{\kilo\gram}$ & \TODO{} & \TODO{} & \TODO{} \\ \hline
        Khuếch đại vi sai, $\Vout$  & $m = \qty{1}{\kilo\gram}$ & \TODO{} & \TODO{} & \TODO{} \\ \hline
        Chuyển đổi V-I, $I$         & $m = \qty{0}{\kilo\gram}$ & \qty{4}{\milli\ampere}  & \TODO{} & \TODO{} \\ \hline
        Chuyển đổi V-I, $I$         & $m = \qty{1}{\kilo\gram}$ & \qty{20}{\milli\ampere} & \TODO{} & \TODO{} \\ \hline
        Đọc dòng, $\Vout$           & $I = \qty{20}{\milli\ampere}$ & \TODO{} & \TODO{} & \TODO{} \\ \hline
    \end{tabularx}
\end{table}

\nhanxet{\TODO{Đánh giá chung các khối trong chương.}}
